% !TeX root=main.tex
\chapter{مروری بر مطالعات انجام شده}
\section{مقدمه}

در این فصل، به بررسی مطالعات و تحقیقات مرتبط با مقیاس‌پذیری بلاکچین می‌پردازیم. ابتدا بیتکوین و معماری آن را بررسی می‌کنیم، سپس راه‌حل‌های مختلف مقیاس‌پذیری را تحلیل می‌کنیم و در نهایت به مقایسه و ارزیابی این راه‌حل‌ها می‌پردازیم.

\section{بیتکوین و معماری آن}

بیتکوین که در سال 2008 توسط ساتوشی ناکاموتو معرفی شد، اولین پیاده‌سازی موفق از یک سیستم پول دیجیتال غیرمتمرکز است \cite{nakamoto2008bitcoin}. بیتکوین از یک معماری تک‌زنجیره‌ای استفاده می‌کند که در آن تمام تراکنش‌ها در یک زنجیره واحد ثبت می‌شوند.

\subsection{الگوریتم اجماع اثبات کار}

بیتکوین از الگوریتم اجماع اثبات کار \lr{(Proof of Work)} استفاده می‌کند. در این الگوریتم، ماینرها برای حل یک مسئله رمزنگاری رقابت می‌کنند و اولین ماینری که راه‌حل را پیدا کند، حق ایجاد بلوک جدید را دارد. این فرآیند نیاز به محاسبات قابل توجهی دارد و به طور متوسط هر 10 دقیقه یک بلوک جدید ایجاد می‌شود.

محدودیت اصلی بیتکوین در اندازه بلوک (حداکثر ۱ مگابایت که بعداً در SegWit به ۴ مگابایت افزایش یافت) و زمان تولید بلوک (۱۰ دقیقه) نهفته است که نتیجه آن حداکثر ۷ تراکنش در ثانیه است. این محدودیت‌ها باعث می‌شود بیتکوین در دوره‌های تقاضای بالا با مشکل تأخیر و هزینه بالا مواجه شود.

\section{راه‌حل‌های لایه دوم}

\subsection{شبکه لایتنینگ}

شبکه لایتنینگ \lr{(Lightning Network)} که در سال 2016 توسط جوزف پون و تادریوس درایجا معرفی شد، یک راه‌حل لایه دوم برای مقیاس‌پذیری بیتکوین است \cite{poon2016lightning}. این شبکه کانال‌های پرداخت خارج از زنجیره را بین کاربران ایجاد می‌کند.

این شبکه قادر به پردازش هزاران تراکنش در ثانیه، با هزینه بسیار پایین و تأیید فوری است؛ با این حال نیاز به باز کردن و مدیریت کانال‌های پرداخت، نقدینگی در کانال‌ها، و پیچیدگی فنی برای کاربران عادی دارد و در صورت مدیریت نادرست ریسک تمرکز ایجاد می‌کند. راه‌حل‌های مشابه دیگری مانند Plasma و Raiden Network برای اتریوم و کانال‌های حالت عمومی \lr{(State Channels)} نیز وجود دارند.

\section{شاردینگ}

شاردینگ یک تکنیک پایگاه داده است که در بلاکچین برای تقسیم شبکه به بخش‌های کوچکتر (شارد) استفاده می‌شود. هر شارد می‌تواند به طور مستقل تراکنش‌ها را پردازش کند.

\subsection{شاردینگ در اتریوم 2.0}

اتریوم ۲/۰ از شاردینگ برای بهبود مقیاس‌پذیری استفاده می‌کند \cite{buterin2014ethereum}. در این معماری شبکه به ۶۴ شارد تقسیم می‌شود، هر شارد مستقل تراکنش‌ها را پردازش می‌کند و ارتباط بین شاردها از طریق beacon chain انجام می‌شود. چالش‌های شاردینگ شامل نیاز به پروتکل‌های پیچیده برای ارتباط بین شاردها، امنیت کمتر هر شارد نسبت به کل شبکه و نیاز به تغییرات اساسی در پروتکل است.
